% tbsrc/pipeline.tex — The Reconstruction Pipeline
% Status: written (first draft). The skeleton common to every worked example:
% a calibration phase (events, mult, contamination, sensitivity) computed once,
% then a reconstruction phase that iterates mlem/osem and is monitored/stopped.
\section{The Reconstruction Pipeline}
\label{sec:pipeline}

The previous sections built every piece of a reconstruction; this one assembles
them into the workflow the worked examples all share. Underneath the different
phantoms and physics, each example is the same two-phase pipeline: a
\emph{calibration} phase that prepares the model inputs once, and a
\emph{reconstruction} phase that iterates. Naming this skeleton is what lets the
examples differ in their data while reusing identical reconstruction code.

\subsection{Calibrate once, iterate many times}
\label{sec:pipeline:twophase}

The objective \eqref{eq:st-listmode} and the MLEM update \eqref{eq:ml-em} read
from four inputs: the event geometry (LOR endpoints), the multiplicative factor
$\code{mult}=\nrm\odot\att$, the additive background
$\code{contamination}=\sct+\rnd$, and the sensitivity image
$\sensv=\Aadj(\code{mult})$. The decisive observation is that \emph{none of these
depends on the activity estimate} $\xv_k$: they are fixed by the scanner and the
object, not by the iterate. So they are computed once, before iterating, and
reused at every step --- and across solvers and datasets, exactly as the OSEM
example shares one $\sensv$ between MLEM and OSEM and between its sharp and
smeared acts. Only the predicted counts $\pred_k$ are recomputed each iteration,
through one forward projection of the current image.

\subsection{The calibration phase}
\label{sec:pipeline:calib}

Calibration turns the geometry and the physical maps into the four inputs. The
events are the listmode LOR endpoints --- read from the scanner for real data, or
sampled on the cylinder for a simulation (Section~\ref{sec:geometry}).
Attenuation comes from projecting the $\mu$-map, $\att=\exp(-\Aop\bm\mu)$
(Section~\ref{sec:datamodel}); normalization $\nrm$ and the background
$\sct+\rnd$ come from external calibration or models. The sensitivity image is
then the back-projection of $\code{mult}$ over the geometric LORs, computed once:
\begin{lstlisting}
xs, xe = sample_lors(scanner, nlors; rng = rng)     # or: read from acquisition
a      = exp.(-joseph3d_fwd(xs, xe, mumap, org, vs))     # attenuation a = exp(-A mu)
mult   = n .* a                                          # n = 1 for the examples
sens   = sensitivity_image(xs, xe, shape, org, vs; weights = mult)   # A^T(mult), once
\end{lstlisting}
In a simulation this phase also generates the data: project the (optionally
blurred) activity, scale to the desired count level, and draw Poisson counts,
\begin{lstlisting}
ybar   = mult .* joseph3d_fwd(xs, xe, gaussian_blur(x_true, fwhm, vs), org, vs) .+ s_plus_r
counts = [rand(rng, Poisson(c)) for c in ybar]
\end{lstlisting}
with $\Gop=\mathbf I$ (no blur) when resolution is switched off. For real data the
counts are simply the recorded events, and everything downstream is identical.

\subsection{The reconstruction phase}
\label{sec:pipeline:recon}

With the inputs prepared, reconstruction assembles the Poisson model and iterates.
The model carries the events, counts, $\code{mult}$, $\code{contamination}$ and
$\sensv$; MLEM (or OSEM, building subset models) runs from a positive image on the
sensitivity support:
\begin{lstlisting}
model = ListmodePoissonModel(xs, xe, sens; img_origin = org, voxsize = vs,
                             counts = counts, mult = mult, contamination = s_plus_r)
x0  = Float32.(sens .> 0)
rec = mlem(model, x0; niter = niter)                       # or osem(subset_models(...), x0)
\end{lstlisting}
This is the whole reconstruction. The physics lives entirely in the prepared
inputs; the solver is generic, which is why switching on attenuation, resolution,
or a background changes only the calibration phase, never these lines.

\subsection{Monitoring and stopping}
\label{sec:pipeline:monitor}

Two quantities track the run. The negative log-likelihood
\code{neg\_log\_likelihood} must decrease monotonically under MLEM
(Section~\ref{sec:mlem}); a non-decrease signals a bug or numerical trouble, so it
is the basic sanity check, available on real data since it needs no truth. When a
reference image \emph{is} known --- the truth or the recoverable $\Gop\xv$ in a
simulation --- the error against it is the quality metric, and it is what reveals
semi-convergence: it falls, reaches a minimum, and rises as noise is fitted
(Section~\ref{sec:acceleration}). The iteration count is therefore not run to
convergence but chosen by early stopping, a reconstruction parameter tuned per
study. The worked examples instrument exactly these quantities --- monotone
likelihood, error against the recoverable target, and the central-slice image ---
and report them alongside each reconstruction.
